% !TeX root = ../papers/nc-p3.tex
\section{二つの二次元次数と四次元の片側対称族}
\label{sec:klein-two-doublets}
前節の残る次数次元$(2,2,1,1)$を扱う。次数$0,1$を二次元とし、
$A=\{0,1\}, B=\{2,3\}$、$\bar i=i+1$と書く（加法は四元群内）。
$b_{ij},c_{ij}$は前節と同じである。異なる組を結ぶときは一次元の
値を双対基底のスカラーで表す。

\begin{theorem}\label{thm:klein-two-doublets-generic}
この型には、幾何学的一般点が反対階数10の標準Koszul AS正則窓となる既約な族がある。
添字保持同型に関する像の次元は4である。一般の偶点スキームは被約な二十点で、
どの三点も共線でない。従ってこの族は、点の印も群作用の印も忘れて三軸族と区別できる。
\end{theorem}

\paragraph{二つの行列の共通非対角成分。}
必要な座標評価階数3以上を課す。同じ組の$i,j$で$b_{ij}=0$なら、
他方の組の両方の$k$について$b_{kj}\ne0$である。
$l=i+j+k$として偶交換を使うと、一次元での対合から$c_{il}=0$を二つ得て矛盾する。
$c$側も同じである。異なる組の$i,j$には
$b_{ij}c_{\bar i,\bar j}=b_{\bar i,j}c_{i,\bar j}$が成り立つ。
二つの列ベクトルはいずれも非零なので、$g_{\bar j}$の倍率を変えれば
$c_{i,\bar j}=b_{ij}$とできる。個別の零成分もこれで含まれる。
以下、薄い次数の二方向を行列$U=(b_{i,2+k})$, $V=(b_{2+k,i})_{k,i}$で表す。
$U$の行順・列順は$i,k=0,1$である。

隣接切断の厚い次数$h=0,1$の行列を
$M_h(i,k)=w_{i,i+h,k,k+h}$とする。偶交換は$M_0$と$M_1^{\mathsf t}$の
対角以外の成分が一致することを与える。同じ組の異なる$i,k$では偶交換そのものであり、
異なる組では双方が同じ薄い二係数の積になる。従って二行列は
\begin{equation}\label{eq:doublet-completions}
 M_0=\begin{pmatrix}S&X\\Y&T\end{pmatrix},\qquad
 M_1^{\mathsf t}=\begin{pmatrix}S'&X\\Y&T'\end{pmatrix},\qquad
 X_{ik}=V_{ki}U_{i,\bar k},\quad Y_{ki}=U_{ik}V_{k,\bar i}.
\end{equation}
局所添字は$0,1$である。$S,S'$と$T,T'$はそれぞれ非対角成分が等しい。
逆にこれと薄いブロックの積表示で全ての偶交換式を尽くす。
実際、二つの偶添字が同じなら恒等式であり、同じ組の相異なる二つなら厚い側の
行列の比較または薄いスカラー積の交換となる。異なる組なら厚い側の成分が薄い側の積に決まる。

$\det Y=-\det X$である。$\det X\ne0$なら階数2条件は
$S,S'$が可逆で$T=YS^{-1}X$, $T'=Y(S')^{-1}X$であることと同値である。
$S'_{00}=i,S'_{11}=j,t=\det S'$を未知とすると、
$Y\operatorname{adj}(S')X-tT$の二つの非対角成分は$i,j,t$の一次式である。
$i,j$に関して解ける開集合では$i,j$は$t$の一次式となり、
行列式条件は高々二次で、既知根$\det S$を持つ。真に二次で残る根も非零なら、
もう一つの解は有理式で一意に得られる。二根が等しい場合もこの有理式は正則であり得る。
$S'=S$なら$g_{\bar j}$と$f_j$の同定で奇交換が成立する。
実際、偶交換も用いると奇交換は厚い二行列が全成分で一致する条件となる。

\paragraph{正則な合流点。}
$U=\left(\begin{smallmatrix}1&2\\2&1\end{smallmatrix}\right)$,
$V=\left(\begin{smallmatrix}3&5\\5&3\end{smallmatrix}\right)$,
$S=\left(\begin{smallmatrix}2&3\\44/87&7\end{smallmatrix}\right)$とする。
二本の一次式は独立で二次式には重根$\det S$がある。
この点の両対称な反対切断は階数8で、両側の核の鉛筆は対角である。
たとえば偶側の二行は平方順に
\[
 (25029,606,-3620,0),\quad (29189,3296,0,-5430)
\]
であり、奇側では$(254616,42891,-392770,0)$と
$(346834,272484,0,-589155)$である。どの二列の小行列式も非零なので両完全交叉は滑らかである。

さらに$U,V$が全成分非零で$\det X\ne0$のとき、第一窓の関係空間は純テンソルを含まない。
関係を四行四列にし、薄い二関係の係数を$z_2,z_3$とすると、
組をまたぐ二つの二行列の行列式は
$U_{00}U_{11}z_2^2-U_{01}U_{10}z_3^2$と
$V_{00}V_{11}z_2^2-V_{01}V_{10}z_3^2$である。
係数行列は$\det X\ne0$によって可逆で、純テンソルなら$z_2=z_3=0$となる。
組内二ブロックの片方は零でなければならないが、
式\eqref{eq:doublet-completions}の両方の列空間が各組へ同型に射影するので他方も零となる。
従って合流点に定理\ref{thm:smooth-apolar-complete}を適用できる。

$U,V,S$の十二係数から上の第二根で作る族は既約で、この合流点を含む。
隣接階数は6、一文字切断は合流点の近くで全階数である。三次・四次の交叉には
テンソルとその縮約から下界があり、上半連続性により合流点の近くで階数は一定となる。
五次交叉は零である開条件を保つ。この有限Frobenius双対を使い、
節\ref{sec:reversible-component}の内部次数別のbar複体の消滅により、
第二根の幾何学的一般点で全次数のKoszul性と左右AS性を得る。
これだけで個々の有理閉点の全次数正則性を判定したことにはならない。

\paragraph{二十点、直線と同型類の次元。}
次の有理入力で反対階数は10である。
\[
 U=\begin{pmatrix}1&2\\3&5\end{pmatrix},\quad
 V=\begin{pmatrix}2&3\\5&7\end{pmatrix},\quad
 S=\begin{pmatrix}2&1\\5&3\end{pmatrix},\quad
 S'=\begin{pmatrix}123/196&1\\5&1335/176\end{pmatrix}.
\]
素数12619で第一評価行列の十五個の最大小行列式を用いる。最初の非零座標を1とした
四つの層で辞書式消去を行うと、幾何点が全て素体上にあり、相異なる二十点を得る。
全ての三点組で三行三列の小行列式の少なくとも一つが非零である。
\path{computations/verify_klein_two_doublets.py}が、上の二入力とこの消去とを検算する。
有限で期待余次元3なのでEagon--Northcott解消から長さは20であり、特殊ファイバーも被約である。
射影性から有限性はこの点の近くで開き、同じ解消と長さにより有限平坦、さらに有限
エタールとなる。その相対三重積のうち相異なる三点の部分は開かつ閉であり、
共線条件の像は閉である。よって標数零の幾何学的一般点にも三点直線はない。
三軸族では有限な場合、各座標面の四符号点の結線が座標頂点を通る。

最後に座標変換を数える。固定座標で薄い二ブロックのうち一方は
二つの四係数ベクトルの外積であり、他方も同じ二ベクトルから決まる。
従って固定テンソルの表示で残る自由度は、二つの薄い次数の値を互いに逆の倍率で
変える一次元だけである。$S$はテンソル成分から直接読める。
射影化も除くと、第二根の射影テンソルの像は$12-1-1=10$次元である。
印付き文字の基底変更は対角で、反対階数10から偶二位置の変換は射影的に一致する。
実際、その反対像は対称二乗全体で、二つの変換による全平方の像が再び対称でなければならない。
奇二位置も$g_{\bar j},f_j$について射影的に同じ倍率になる。
これは組をまたぐ二つの奇添字では、規格化したテンソルが奇交換で等しく、
その成分が全て非零だからである。逆に偶・奇でそれぞれ独立に同じ対角変換を施してよく、
有効な射影対角群は六次元である。

一般の正則点では両奇点スキームも偶側と同型で有限被約である。
三つの射影表示で二次式は点全体で消え得ない（最大小行列式の最小次数は4）。
従ってこれらの有限集合の射影安定群は有限である。実際、各点を固定する非スカラー行列の
固有直線全体は二つの超平面の和集合に入る。よってテンソルを保つ射影変換群も有限で、
四元群作用の選び直しは有限通りである。各文字が重複しないので、作用を固定した
基底の選び直しは上で数えた対角変換と有限置換で尽きる。
定理\ref{thm:frobenius-tensor-reconstruction}によって添字保持同型を取った像は
$10-6=4$次元となる。
